%=============================================================================
% Chapter 12: Machine Learning Pipeline
%=============================================================================
\chapter{Machine Learning Pipeline}

\section{Complete Pipeline Overview}

The FloodEye ML pipeline spans two environments: the training pipeline runs in Python (Google Colab), and the inference pipeline runs in TypeScript (browser). Figure~\ref{fig:ml_pipeline} illustrates the complete end-to-end pipeline.

\begin{figure}[htbp]
\centering
\resizebox{!}{0.85\textheight}{%
\begin{tikzpicture}[node distance=0.4cm]
  % Training pipeline
  \node[startend, text width=3cm] (dataset) {Dataset\\\texttt{.\allowbreak{}csv}};
  \node[process, text width=3cm, below=of dataset] (eda) {EDA \&\\Visualisation};
  \node[process, text width=3cm, below=of eda] (clean) {Data Cleaning\\(Sort, Index)};
  \node[process, text width=3cm, below=of clean] (feateng) {Feature\\Engineering};
  \node[process, text width=3cm, below=of feateng] (split) {Chronological\\Split (70/15/15)};
  \node[process, text width=3cm, below=of split] (scale) {StandardScaler\\Normalisation};
  \node[process, text width=3cm, below=of scale] (windows) {Sliding Window\\(LOOKBACK=48)};
  \node[process, text width=3cm, below=of windows] (build) {Build Model\\(CNN+BiLSTM+Att)};
  \node[process, text width=3cm, below=of build] (train) {Train\\(Adam, MSE)};
  \node[process, text width=3cm, below=of train] (validate) {EarlyStopping\\Validation};
  \node[process, text width=3cm, below=of validate] (eval) {Evaluation\\(MAE, Plots)};
  \node[process, text width=3cm, below=of eval] (serial) {Serialise\\(\texttt{.\allowbreak{}h5})};
  \node[process, text width=3cm, below=of serial] (convert) {TFJS Convert\\+ Patch Script};
  \node[startend, text width=3cm, below=of convert] (artifacts) {Export\\Artifacts};

  % Inference pipeline (right column)
  \node[process, text width=3cm, right=3.5cm of build] (load) {Load Model\\(browser)};
  \node[process, text width=3cm, below=of load] (live) {Live History\\(last 48 readings)};
  \node[process, text width=3cm, below=of live] (featts) {Feature\\Engineering (TS)};
  \node[process, text width=3cm, below=of featts] (norm) {Normalise\\(scalers.json)};
  \node[process, text width=3cm, below=of norm] (infer) {TF.js\\Inference};
  \node[process, text width=3cm, below=of infer] (denorm) {Denormalise\\Prediction};
  \node[startend, text width=3cm, below=of denorm] (output) {Risk Level \&\\4h Forecast};

  % Training arrows
  \draw[arrow] (dataset) -- (eda);
  \draw[arrow] (eda) -- (clean);
  \draw[arrow] (clean) -- (feateng);
  \draw[arrow] (feateng) -- (split);
  \draw[arrow] (split) -- (scale);
  \draw[arrow] (scale) -- (windows);
  \draw[arrow] (windows) -- (build);
  \draw[arrow] (build) -- (train);
  \draw[arrow] (train) -- (validate);
  \draw[arrow] (validate) -- (eval);
  \draw[arrow] (eval) -- (serial);
  \draw[arrow] (serial) -- (convert);
  \draw[arrow] (convert) -- (artifacts);

  % Inference arrows
  \draw[arrow] (load) -- (live);
  \draw[arrow] (live) -- (featts);
  \draw[arrow] (featts) -- (norm);
  \draw[arrow] (norm) -- (infer);
  \draw[arrow] (infer) -- (denorm);
  \draw[arrow] (denorm) -- (output);

  % Cross-pipeline connection
  \draw[arrowg,dashed] (artifacts) -| node[above,font=\tiny,pos=0.3]{deploy to /public} (load);

\end{tikzpicture}%
}
\caption{Complete ML Training and Inference Pipeline}
\label{fig:ml_pipeline}
\end{figure}

\section{Stage-by-Stage Explanation}

\subsection{Stage 1: Dataset}
The raw CSV dataset (\texttt{assam\_\allowbreak{}flood\_\allowbreak{}sensor\_\allowbreak{}data.\allowbreak{}csv}) contains timestamped multivariate sensor readings with a look-ahead water level target. The dataset is loaded into a Pandas DataFrame and sorted chronologically.

\subsection{Stage 2: EDA and Visualisation}
Seaborn and Matplotlib are used to generate the four-panel EDA plot (\texttt{flood\_\allowbreak{}plots/\allowbreak{}eda\_\allowbreak{}plots.\allowbreak{}png}), providing insight into the temporal distribution, seasonal patterns, feature correlations, and key relationships.

\subsection{Stage 3: Data Cleaning}
The datetime index is constructed, the DataFrame is sorted, and any null values are handled. The data is confirmed to be regularly sampled and chronologically ordered.

\subsection{Stage 4: Feature Engineering}
The \texttt{prepare\_\allowbreak{}sequences()} function engineers 11 features from the raw columns: four cyclical time encodings, five raw sensor readings, and two 3-hour trend features.

\subsection{Stage 5: Chronological Split}
The dataset is split 70/15/15 without shuffling, preserving temporal order to prevent data leakage.

\subsection{Stage 6: StandardScaler Normalisation}
Scalers are fitted on the training split only. Both X (features) and Y (target) are normalised to zero mean and unit variance. Scaler parameters are exported to \texttt{scalers.\allowbreak{}json}.

\subsection{Stage 7: Sliding Window Creation}
Overlapping windows of length 48 are created over the scaled feature matrix, producing input tensors \texttt{(N,\allowbreak{} 48,\allowbreak{} 11)} and target scalars \texttt{(N,\allowbreak{} 1)}.

\subsection{Stage 8: Build Model}
The 1D-CNN + Bidirectional LSTM + Attention architecture is constructed using the Keras functional API. The custom \texttt{AttentionLayer} is defined as a Keras Layer subclass.

\subsection{Stage 9: Training}
Model is compiled with the Adam optimizer and MSE loss. Training runs for up to 50 epochs with batch size 32.

\subsection{Stage 10: EarlyStopping Validation}
Training is halted when validation loss fails to improve for 10 consecutive epochs; the best weights are restored. ReduceLROnPlateau halves the learning rate when validation loss plateaus for 5 epochs.

\subsection{Stage 11: Evaluation}
Four evaluation plots are generated (\texttt{flood\_\allowbreak{}plots/\allowbreak{}evaluation\_\allowbreak{}plots.\allowbreak{}png}): actual vs predicted on the full test set, zoomed peak event, residuals, and attention weights.

\subsection{Stage 12: Serialisation}
The trained model is saved as \texttt{flood\_\allowbreak{}model.\allowbreak{}h5}. Scaler parameters (\texttt{scalers.\allowbreak{}json}), model configuration (\texttt{config.\allowbreak{}json}), and sample data (\texttt{sample\_\allowbreak{}data.\allowbreak{}json}) are saved.

\subsection{Stage 13: TFJS Conversion and Patching}
The \texttt{tensorflowjs\_\allowbreak{}converter} converts \texttt{flood\_\allowbreak{}model.\allowbreak{}h5} to LayersModel format. The \texttt{patch\_\allowbreak{}model.\allowbreak{}js} script patches the resulting \texttt{model.\allowbreak{}json} for Keras 2 compatibility, fixes weight name paths, and replaces the Orthogonal initialiser.

\subsection{Stage 14: Export Artifacts}
All artifacts are packaged: \texttt{tfjs\_\allowbreak{}model/\allowbreak{}model.\allowbreak{}json}, \texttt{tfjs\_\allowbreak{}model/\allowbreak{}group1-shard1of1.\allowbreak{}bin}, \texttt{scalers.\allowbreak{}json}, \texttt{config.\allowbreak{}json}, \texttt{sample\_\allowbreak{}data.\allowbreak{}json}. These are copied to \texttt{public/\allowbreak{}} for browser access.

\subsection{Stage 15: Browser Model Loading}
The \texttt{loadModel()} function in \texttt{lib/\allowbreak{}flood-model.\allowbreak{}ts} fetches all artifacts in parallel, manually loads weights, and registers the custom AttentionLayer.

\subsection{Stage 16--19: Browser Inference}
Live history data from the TelemetryProvider is feature-engineered, normalised, shaped into a 3D tensor, passed through the model, denormalised, and mapped to a risk level.

\section{Key Pipeline Design Decisions}

\begin{table}[H]
\centering
\caption{Key ML Pipeline Design Decisions}
\begin{tabularx}{\linewidth}{|l|X|}
\hline
\rowcolor{FloodLight}
\textbf{Decision} & \textbf{Rationale} \\
\hline
Chronological split & Prevents temporal data leakage; reflects real-world evaluation \\
48-step lookback & Captures 48 hours of context; sufficient for monsoon trend detection \\
BiLSTM over LSTM & Forward+backward context improves prediction at any position in the sequence \\
Attention layer & Interpretability + focus on critical time steps (pressure drops, rapid rises) \\
Client-side inference & Zero server cost for ML; scales to unlimited concurrent users \\
Manual weight loading & Bypasses TFJS naming issues caused by Keras 3 export format \\
Sample data padding & Allows inference to begin immediately without 48h of live data \\
\hline
\end{tabularx}
\end{table}
